\section{Lesson 5: The Compactness--Proportionality Pareto Frontier Is Steep}
\label{sec:lesson5}

The fifth and most technically precise lesson from the E-track experiments is
that the Pareto frontier between compactness and proportionality is steep: a
one percentage point improvement in seat--vote proportionality costs
approximately 0.015 units of mean Polsby--Popper, roughly 4\% of the baseline
value.  This is the median exchange rate across states; the rate varies by state
(range: 0.011--0.023~PP/pp), reflecting differences in urban concentration.  We
report the median value 0.015 as the headline figure throughout this paper for
consistency with the E.0 cross-system synthesis \citep{dellaLibera2026e0}.
No feasible single-election geographic system escapes this trade-off.

\subsection{Deriving the Exchange Rate}

Paper E.5 implements a parametric optimization that trades compactness for
proportionality by varying the weight on the efficiency-gap criterion in
the METIS objective function \citep{dellaLibera2026e5}.  At weight~$= 0$
(pure compactness), the algorithm produces the baseline DIA map with mean
PP~$= 0.367$ and seat--vote deviation~$= 4.2$~pp.  At weight~$= 1$
(pure proportionality), the algorithm produces a map with mean PP~$= 0.294$
and seat--vote deviation~$= 0.3$~pp.  The relationship is approximately linear
across the weight range, giving a \emph{within-system} exchange rate of:
\[
  \frac{\Delta \mathrm{PP}}{\Delta \text{deviation}} \approx
  \frac{0.367 - 0.294}{4.2 - 0.3} = \frac{0.073}{3.9} \approx 0.019 \text{ PP per pp}.
\]
This within-system rate (0.019) measures the slope of the E.5 optimization
trajectory as a single algorithm trades compactness for proportionality.  The
E.0 synthesis \citep{dellaLibera2026e0} estimates a \emph{cross-system} rate
of 0.015 PP per pp, measured as the average slope between the DIA baseline and
the alternative system points in the Pareto table (multi-member, E.5, etc.).
These two numbers measure different phenomena.  The cross-system rate (0.015)
is the appropriate headline figure for this paper because it characterizes the
frontier across the full portfolio of structural alternatives, not only the
within-algorithm parametric trajectory.  We therefore use 0.015 throughout.

\subsection{What the Steep Frontier Means}

An exchange rate of 0.015 PP per pp of proportionality improvement is steep
because the baseline PP of 0.367 is already substantially higher than most
real-world maps.  The practical implication is:

\begin{itemize}
  \item Moving from 4.2~pp seat--vote deviation (baseline) to 2.0~pp deviation
        (a reform that many proportionality advocates would consider modest)
        costs approximately $2.2 \times 0.015 = 0.033$ PP, reducing mean PP
        from 0.367 to 0.334 --- a meaningful loss of compactness.
  \item Achieving 0.5~pp deviation (near-perfect proportionality) costs
        approximately $3.7 \times 0.015 = 0.056$ PP, reducing mean PP
        to 0.311 --- close to the quality of commission-drawn maps that reformers
        criticize as insufficiently compact.
\end{itemize}

In other words, algorithmic proportionality reform would, at the level required
to achieve near-perfect seat--vote symmetry, produce maps no more compact than
the commission-drawn maps that algorithmic redistricting is intended to improve
upon.  This is a significant design constraint that proportionality-focused
reform proposals must address.

\subsection{Why No System Escapes the Frontier}

The steep Pareto frontier is a consequence of geographic sorting: the same
geographic sorting that produces the seat--vote gap also makes it geometrically
expensive to fix.  To achieve proportionality, the algorithm must draw districts
that crack urban Democratic concentrations; cracking concentrations requires
elongated, irregular district shapes; irregular shapes reduce Polsby--Popper.
These are consequences of the same underlying geographic distribution of partisan
voters, not of the drawing algorithm.  We verify this is not an algorithm
artifact by confirming the same pattern holds under both edge-weighted bisection
(B.2) and standard bisection (B.1) configurations: the exchange rate under
standard bisection (0.016~PP/pp) is statistically indistinguishable from the
0.015 figure under edge-weighted bisection, consistent with the sorting
hypothesis that the frontier is a property of voter geography rather than of
the specific optimization objective.

This can be demonstrated by a thought experiment.  Suppose we had a perfectly
unbiased redistricting algorithm (if such a thing were possible) that drew
districts from a random partition of the census-tract graph subject only to
population equality and contiguity.  Such an algorithm would produce maps with
mean PP near 0.15--0.20 (the compactness of random planar graph partitions)
and seat--vote deviations near zero (because partisan sorting would be broken
up randomly).  Both the compactness and the proportionality of the baseline
DIA are improvements over this random baseline, but they are improvements on
different objectives, and improving further on one requires accepting deterioration
on the other.

The key insight is that the sorting constraint is not a feature that can be
``engineered out'' of a geographic district system by choosing a different
algorithm.  It is a property of the data --- of where voters live --- that
any geographically coherent drawing must contend with.

\subsection{Implications for Reform Evaluation}

The steepness of the Pareto frontier implies that reform proposals should be
evaluated not just on the proportionality they achieve but on the compactness
they sacrifice.  A proposal that achieves 2~pp seat--vote deviation at a cost
of 0.033 PP may be acceptable; a proposal that achieves 0.5~pp deviation at
a cost of 0.056 PP may not be.

This framework suggests a useful reformulation of the redistricting criteria
debate.  Rather than treating compactness and proportionality as competing
criteria whose relative weights are determined by legislative or judicial
preference, we can treat the Pareto frontier itself as an objective criterion:
\emph{any redistricting plan should lie on or near the Pareto frontier for its
state's geography.}  A plan that is less proportional \emph{and} less compact
than what the frontier permits is objectively inferior.  The debate about where
on the frontier to locate the plan --- how to trade compactness for
proportionality --- is a normative question, but it should be conducted with
full knowledge of what the trade-off actually costs.

The E-track research program provides the empirical foundation for this
frontier-based evaluation.  The exchange rates documented in E.5 and
corroborated across E.1--E.5 give reformers and courts a quantitative basis
for assessing whether a proposed redistricting plan is making a legitimate
normative trade-off or simply producing a map that is both less compact and
less proportional than the state's geography permits.
